\documentclass[10pt, a4paper]{article}

% ── PACKAGES ──────────────────────────────────────────────────
\usepackage[utf8]{inputenc}
\usepackage[T1]{fontenc}
\usepackage{geometry}
\geometry{top=2.5cm,bottom=2.5cm,
          left=2.5cm,right=2.5cm}
\usepackage{hyperref}
\hypersetup{
  colorlinks=true,urlcolor=blue,
  linkcolor=blue,citecolor=blue,
  pdftitle={Sleep as Attractor Basin
    Navigation Cost: A Unified Geometric
    Theory of Sleep Function, Topology,
    and Deep-Time Survival},
  pdfauthor={Eric Robert Lawson /
    OrganismCore}
}
\usepackage{microtype,parskip,booktabs}
\usepackage{array,tabularx,longtable}
\usepackage{xcolor,mdframed,enumitem}
\usepackage{titlesec,lmodern,fancyhdr}
\usepackage{lastpage,amsmath,amssymb}
\usepackage{abstract}

% ── COLOURS ────���──────────────────────────────────────────────
\definecolor{headergreen}{RGB}{20,100,60}
\definecolor{rulecolor}{RGB}{20,100,60}
\definecolor{claimbox}{RGB}{20,100,60}
\definecolor{shadelight}{RGB}{248,252,248}

% ── HEADER / FOOTER ───────────────────────────────────────────
\pagestyle{fancy}\fancyhf{}
\renewcommand{\headrulewidth}{0.4pt}
\fancyhead[L]{\small\color{headergreen}
  Lawson / OrganismCore}
\fancyhead[R]{\small\color{headergreen}
  Sleep as Attractor Basin Navigation
  Cost --- 2026-03-18}
\fancyfoot[C]{\small
  Page \thepage\ of \pageref{LastPage}}

% ── SECTION FORMAT ────────────────────────────────────────────
\titleformat{\section}
  {\large\bfseries\color{headergreen}}
  {\thesection.}{0.5em}{}[
  \vspace{-0.3em}\textcolor{rulecolor}
  {\rule{\linewidth}{0.6pt}}]
\titleformat{\subsection}
  {\normalsize\bfseries\color{headergreen}}
  {\thesubsection.}{0.5em}{}
\titleformat{\subsubsection}
  {\small\bfseries}
  {\thesubsubsection.}{0.5em}{}
\setlength{\parskip}{0.5em}
\setlength{\parindent}{0pt}

% ══════════════════════════════════════════════════════════════
\begin{document}

% ── TITLE BLOCK ───────────────────────────────────────────────
\begin{center}
  {\LARGE\bfseries\color{headergreen}
  Sleep as Attractor Basin Navigation
  Cost:\\[0.3em]
  A Unified Geometric Theory of Sleep
  Function, Topology,\\[0.3em]
  and Deep-Time Survival}\\[0.8em]
  {\normalsize
  \textbf{Eric Robert Lawson}\\
  OrganismCore (Independent Research)\\
  \href{mailto:OrganismCore@proton.me}
  {OrganismCore@proton.me}\\
  ORCID:
  \href{https://orcid.org/0009-0002-0414-6544}
  {0009-0002-0414-6544}\\[0.4em]
  \textit{Manuscript deposited as
  reasoning artifact series,
  OrganismCore, 2026-03-18.}\\
  \textit{Repository:}
  \href{https://github.com/Eric-Robert-Lawson/attractor-oncology}
  {\texttt{github.com/Eric-Robert-Lawson/
  attractor-oncology}}\\
  \textit{Theory paper:}
  \href{https://doi.org/10.5281/zenodo.19094935}
  {\texttt{doi:10.5281/zenodo.19094935}}
  }
\end{center}

\vspace{0.5em}
\noindent\textcolor{rulecolor}
{\rule{\linewidth}{1pt}}
\vspace{0.5em}

% ── ABSTRACT ──────────────────────────────────────────────────
\begin{mdframed}[backgroundcolor=shadelight,
  linecolor=headergreen,linewidth=0.8pt,
  innertopmargin=0.7em,innerbottommargin=0.7em,
  innerleftmargin=0.9em,innerrightmargin=0.9em]
\textbf{Abstract}

\smallskip
We present a unified geometric theory
of sleep derived from the attractor
landscape framework
$S + N + G \rightarrow R$, where $S$
is the structural substrate (genome and
somatic integrity), $N$ is the coherence
gradient field (circadian clock, social
environment, and light), and $G$ is the
global attractor landscape, physically
implemented as the organism's epigenetic
state. The framework makes three
inseparable predictions about rest:
(1) \textit{Quantity} --- the total cost
of occupying a position in the attractor
landscape must be discharged per cycle,
and duration of discharge tracks total
cost; (2) \textit{Topology} --- the
architecture of the discharge process is
itself a geometric prediction of the
landscape position, not an independent
biological variable; and (3)
\textit{Substrate} --- the cost of basin
navigation is fundamentally oxidative,
with reactive oxygen species (ROS)
accumulation as the proximate death
mechanism in sleep deprivation (Vaccaro
et al., 2020), and the glymphatic system
as the topological mechanism of discharge
within consolidated bilateral sleep.

From these three predictions we derive:
eight sleep topologies (T1 through T8);
two deep-time survival strategies (Lock
and Navigator); an unstable zone; a
complete three-channel sonar instrument
(duration, topology, behavioural
flexibility); and the dimensional
invariant --- the oldest surviving
lineages sleep the least and the minimum
of the sleep spectrum is the maximum of
the survival spectrum.

We further identify that (a) the four
dominant sleep function theories in the
current literature (synaptic homeostasis,
energy conservation, glymphatic
clearance, memory consolidation) are
correct descriptions of mechanisms
within a single topology class (T3)
and are subsumed by the framework as
special cases; (b) the derivation chain
from cancer biology (FOXA1/EZH2 ratio;
Lawson, 2024) closes on itself through
BMAL1 --- loss of the master circadian
clock gene makes the cellular repair
basin epigenetically inaccessible,
producing a cancer attractor at the
cellular level geometrically identical
to the organism-level failure it
initiated; (c) social information is
a primary $N$ field co-equal to light,
with dominance rank negatively
correlating with sleep quality
proportional to the social navigational
demand of the dominant position; and
(d) artificial light at night (ALAN)
is a global primary $N$ field disruption
operating simultaneously across
thousands of species.

Eleven a-priori predictions
(P-INVARIANT-1 through P-INVARIANT-11)
and five closure predictions
(P-CLOSURE-1 through P-CLOSURE-5)
are stated. None are derivable from
any existing sleep function theory.
All were stated before literature
confirmation was sought.
\end{mdframed}

\vspace{0.6em}

% ══════════════════════════════════════════════════════════════
\section{Introduction: The Problem With
Four Theories}
\label{sec:intro}
% ══════════════════════════════════════════════════════════════

Sleep is universal across animal life
and has persisted across 600 million
years of evolution despite imposing an
obvious and severe fitness cost: a
sleeping organism cannot feed, reproduce,
or detect predators. The persistence
of sleep therefore implies a benefit
sufficiently large to outweigh this
cost continuously across deep time.

The current literature offers four
competing frameworks for what that
benefit is:

\begin{enumerate}[itemsep=0.3em]
  \item \textbf{Synaptic Homeostasis
    (SHY):} Sleep downscales synaptic
    potentiation accumulated during
    waking, preventing synaptic
    saturation and restoring the
    brain's capacity for new learning
    \cite{tononi2006sleep}.

  \item \textbf{Energy conservation:}
    Sleep reduces metabolic demand
    and allows cellular energy reserve
    replenishment \cite{siegel2005clues}.

  \item \textbf{Glymphatic clearance:}
    Sleep drives cerebrospinal fluid
    pulsation through perivascular
    channels, clearing neurotoxic
    metabolites including amyloid-beta
    \cite{xie2013sleep}.

  \item \textbf{Memory consolidation:}
    Sleep reactivates and stabilises
    memory traces formed during waking
    \cite{stickgold2005sleep}.
\end{enumerate}

The 2023--2024 consensus is explicit:
no single theory explains all aspects
of sleep \cite{frank2024sleep}. The
frameworks are treated as complementary
and overlapping. This theoretical
pluralism is the symptom of a missing
geometric structure.

\textbf{The problem:} None of these
four frameworks generates the following
observations:

\begin{itemize}[itemsep=0.2em]
  \item The oldest surviving lineages
    on Earth sleep the least.
  \item Organisms that sleep almost
    nothing (horseshoe crabs, sponges,
    nautilus) have survived every mass
    extinction for hundreds of millions
    of years.
  \item Organisms that sleep the most
    (bats, koalas, cats) are not the
    most cognitively complex.
  \item The dolphin has dissolved the
    barrier between waking and sleep
    and permanently occupies an
    intermediate state.
  \item The octopus has independently
    evolved a two-phase sleep
    architecture in a brain with
    no cerebral cortex.
  \item A crocodilian sleeps 2--6 hours
    of true bilateral EEG sleep and
    12--16 hours of distributed rest,
    not because it is cognitively
    simple but because its ectothermic
    metabolic architecture determines
    the only topologically feasible
    discharge architecture for
    its navigational demand.
\end{itemize}

These observations are not anomalies.
They are the geometric signal. The
theory that explains them must be
geometric, not mechanistic.

We propose such a theory.

% ══════════════════════════════════════════════════════════════
\section{The Attractor Geometry
Framework}
\label{sec:framework}
% ══════════════════════════════════════════════════════════════

\subsection{The Fundamental Equation}

\[
S + N + G \;\rightarrow\; R
\]

\begin{itemize}[itemsep=0.2em]
  \item $S$ = structural substrate
    (genome, epigenetically accessible
    gene programmes, somatic tissue
    integrity, stem cell reserve)
  \item $N$ = coherence gradient field
    (light, social information, and
    circadian clock genes as the
    molecular implementation of $N$
    at the cellular level)
  \item $G$ = global attractor landscape,
    physically implemented as the
    organism's epigenetic state:
    histone modification patterns,
    DNA methylation, and chromatin
    accessibility across the genome
  \item $R$ = resulting state or
    phenotype
\end{itemize}

An organism's position in $G$ determines
four properties simultaneously:
(1) the total cost of occupying that
position; (2) the survival trajectory;
(3) the topology of the discharge process;
and (4) the biochemical substrate of
the cost.

\subsection{Sleep as Basin-Navigation
Cost Discharge}

We define sleep as the generic process
by which the accumulated cost of
basin occupation in $G$ is discharged.
The cost has three axes:

\begin{description}[itemsep=0.3em]
  \item[Axis~1 --- Between-basin
    switching cost.] The organism
    navigates between multiple
    accessible basins in $G$.
    Each transition requires bilateral
    synaptic reconfiguration and
    generates ROS in proportion to
    the neural firing intensity of
    the transition.

  \item[Axis~2 --- Within-basin
    metabolic cost.] A single
    high-cost basin consumes all
    available metabolic resources.
    Sleep is required for metabolic
    recovery and oxidative debt
    clearance rather than synaptic
    reconfiguration.

  \item[Axis~3 --- Within-basin
    computational intensity.] Extreme
    neural computation within a single
    specialised basin generates high
    ROS from neural firing intensity.
    Discharge is required for both
    synaptic and oxidative cost.
\end{description}

\subsection{The Oxidative Substrate}

The cost of navigation is fundamentally
oxidative. This is confirmed by the
finding that sleep deprivation kills
via reactive oxygen species accumulation
specifically in the intestinal epithelium
of \textit{Drosophila}, and that
antioxidant supplementation prevents
death under total sleep deprivation even
without additional sleep
\cite{vaccaro2020sleep}. The gut is the
tissue most directly responsible for
substrate acquisition driving the
navigation, and its epithelium bears
the highest peripheral oxidative cost.

Sleep discharges two costs simultaneously:
the synaptic cost (via slow-wave
downscaling, the Synaptic Homeostasis
mechanism \cite{tononi2006sleep}) and
the oxidative cost (via glymphatic
pulsation during NREM slow-wave sleep,
driven by interstitial space expansion
and AQP4-mediated CSF flow through
perivascular channels \cite{xie2013sleep}).
The glymphatic system is the physical
mechanism of oxidative cost discharge
within the consolidated bilateral sleep
topology.

\subsection{G is the Epigenetic State}

The attractor landscape $G$ is not
an abstract theoretical object. Its
physical substrate is the epigenetic
state of the organism's cells.

Basin depths in $G$ correspond to the
depth of epigenetic entrenchment of
the corresponding state. Barrier heights
correspond to the chromatin
reconfiguration energy required to
transition between states. Sleep
maintains $G$ daily through
sleep-dependent histone deacetylation
(restoring chromatin accessibility
patterns to their functional state)
and DNA methylation maintenance
\cite{lappalainen2023sleep}.

Chronic sleep deprivation degrades $G$
irreversibly: progressive histone
modification dysregulation, DNA
methylation drift, pleiotrophin depletion
leading to neuronal death (removing
basins permanently), and
NF-$\kappa$B-driven inflammatory
landscape remodelling \cite{acs2023sleep}.
The irreversible components of chronic
sleep deprivation damage are $G$
degradation past the restoration
threshold.

% ══════════════════════════════════════════════════════════════
\section{The Eight Sleep Topologies}
\label{sec:topologies}
% ══════════════════════════════════════════════════════════════

The framework predicts that different
landscape positions produce different
sleep architectures. Each topology is
the only feasible discharge architecture
for its corresponding landscape position
given the full constraint set.

\subsection{T1: Metabolic Quiescence}

Deepest Lock organisms with minimal
neural architecture. Zero navigational
demand. Zero synaptic reconfiguration
cost. Minimal ROS generation. No sleep
in any neurological sense; circadian
biochemical rhythms only.

\textit{Confirmed:} Sponge ($\sim$600\,Myr),
horseshoe crab ($\sim$450\,Myr), most
jellyfish (500--700\,Myr).

\subsection{T2: Distributed Bilateral
Rest}

Ectothermic Navigators. Navigator-level
navigational cost but ROS generation
rate an order of magnitude lower than
endothermy. Energy budget for full
bilateral wakefulness limited per day.

True bilateral EEG sleep for 2--6h
discharging synaptic cost. Extended
distributed bilateral rest for 12--16h
clearing the low-rate oxidative debt
continuously. Total Navigator-equivalent
discharge. Why bilateral and not
unihemispheric: the switching cost is
between ecological basins, bilaterally
symmetric, requiring no hemispheric
asymmetry in discharge.

\textit{Confirmed:} Crocodilian
(bilateral EEG throughout rest periods,
Kelly \& Lesku, 2015, 2019). One-eye-open
vigilance is behavioural (brainstem
alerting), not unihemispheric in the
EEG sense.

\subsection{T3: Consolidated Bilateral
Deep Sleep}

Endothermic organisms. High ROS
generation rate during waking requires
concentrated high-efficiency glymphatic
clearance during sustained NREM.
Bilateral synaptic reconfiguration
discharged during same concentrated
window.

\textit{Confirmed:} All endothermic
organisms. Duration within T3 is a
valid direct proxy for total cost
because metabolic architecture is
held constant.

\subsection{T4: Permanent Basin~Omega ---
Dissolved Barrier}

Open-ocean Navigators where full
bilateral sleep is lethal. The cost
of entering the sleep basin's exit
barrier exceeds the accumulated cost
of remaining partially awake
indefinitely. Barrier between waking
and sleep is dissolved. Hemispheric
alternation in continuous rotation.

\textit{Confirmed:} Dolphin, porpoise,
some seals. Frigatebirds during
multi-day oceanic flight.

\subsection{T5: Fragmented Shallow Sleep
--- External Suppression}

Deep single basin under chronic external
pressure interrupting sleep before
discharge completes. Sleep basin is
behaviourally shallow, not geometrically
shallow. Debt accumulates and is
discharged as rebound when pressure is
removed. Rebound is the definitive
diagnostic.

\textit{Confirmed:} Giraffe.

\subsection{T6: Transition-Cost Sleep ---
Developmental Basin Crossing}

Sleep required not to discharge ongoing
occupation cost but to pay the cost of
a developmental basin transition.
Duration proportional to transition
magnitude, not to ongoing navigational
demand. Cannot be postponed without
developmental failure.

\textit{Confirmed:} \textit{C.~elegans}
lethargus at each larval transition
\cite{allada2008invertebrate}. Broadly,
all animals with discrete developmental
basin transitions.

\subsection{T7: Immune-Basin Sleep ---
Defensive Basin Occupation}

Pathogen invasion creates a new
high-cost basin: the immune response.
Sleep is the discharge process for
immune-basin occupation cost. Cytokines
are the $N$ field signal specifying
immune basin entry. Duration proportional
to immune response intensity.

\textit{Confirmed:}
\textit{C.~elegans} sickness quiescence
\cite{hill2014cellular}. Human and
mammalian fever sleep. Confirmed across
all major animal phyla.

\subsection{T8: Suspended State with
Blocked Discharge}

Hibernation/torpor. Both navigational
activity and sleep architecture are
suspended by temperature constraints.
Homeostatic pressure continues to
build in the complete absence of
navigational activity. Periodic arousals
driven by homeostatic pressure reaching
threshold --- metabolically expensive
warming purely to enable intense
slow-wave sleep discharge. Post-emergence
rebound qualitatively different from T5:
maximum debt with zero navigational
cause.

\textit{Confirmed:} Ground squirrel,
bear, bat, hedgehog. Post-emergence
rebound sleep proportional to
hibernation duration.

\subsection{The Eight-Topology Prediction}

The eight topologies are not a closed
taxonomy. They are the geometrically
derivable solutions given the known
constraint set. Any new topology
discovered empirically will be derivable
from the same principles: the interaction
of navigational cost structure, ROS
generation rate, metabolic architecture,
environmental lethality constraints on
full sleep, developmental transition
requirements, and immune-basin
occupation. The framework is not a
list. It is a geometry.

% ══════════════════════════════════════════════════════════════
\section{The Two Survival Strategies
and the Unstable Zone}
\label{sec:strategies}
% ══════════════════════════════════════════════════════════════

\subsection{The Lock Strategy}

Single basin. Extremely deep. Very high
exit barrier. Minimal metabolism. Minimal
ROS generation. Sleep duration approaching
zero in simplest forms. Topology T1 or T2.
Survival horizon: 300--600\,Myr+.

Convergent across seven independent
phyla: Porifera, Cnidaria, Mollusca,
Chelicerata, Chondrichthyes,
Sarcopterygii, Lepidosauria.

Survival mechanism: basin so deep it
is perturbation-invariant. The Permian--
Triassic event eliminated 96\% of
marine species. The horseshoe crab
(T1, near-zero sleep, 450\,Myr) survived.

\subsection{The Navigator Strategy}

Multiple basins. Shallow barriers.
High switching capacity. Ecological
generalism. Moderate to high sleep
equivalent. Topology T2, T3, or T4
depending on metabolic architecture and
environmental constraints. Survival
horizon: 100--300\,Myr.

Convergent across six+ independent
lineages across T2, T3, and T4.
The critical confirmation: octopus
(Cephalopoda, $\sim$300\,Myr) has
independently evolved a two-phase
sleep (quiet plus active REM-like)
in a brain with no cerebral cortex
\cite{medeiros2021sleeping}. The same
topology from completely different
neural architecture. Topology is
determined by navigational demand,
not by neural architecture.

\subsection{The Unstable Zone}

Shallow single basin. High barriers
to alternatives. Not deep enough for
perturbation-invariance. Not flexible
enough to reach alternatives when
niche disappears. Extinction concentrated
here. Examples: \textit{T.~rex},
mammoth, dodo.

\subsection{The Dimensional Invariant}

\begin{quote}
\itshape
The oldest surviving lineages sleep the
least. The minimum of the sleep spectrum
is the maximum of the survival spectrum.
This is not a coincidence. This is the
geometry.
\end{quote}

The survival--sleep duration relationship
is U-shaped with a flat minimum. The
minimum of sleep duration corresponds to
the maximum survival horizon (Lock
strategy). The moderate range is the
Navigator survival zone. The high-duration
range is uncertain --- survival depends
on niche stability or Basin~$\Omega$
depth.

% ══════════════════════════════════════════════════════════════
\section{The N Field: Clock, Social
Information, and Light}
\label{sec:nfield}
% ══════════════════════════════════════════════════════════════

The coherence gradient field $N$
has three confirmed physical
implementations:

\subsection{The Circadian Molecular Clock}

BMAL1/CLOCK/PER/CRY forms the molecular
$N$ field that specifies when cells and
organisms transition between active and
rest/repair basins. BMAL1/CLOCK acts
as a histone acetyltransferase --- it
directly opens chromatin at repair
gene loci, implementing the transition
into the repair basin at the molecular
level.

BMAL1 loss at the cellular level is
the molecular equivalent of an organism
that cannot enter its sleep basin.
BMAL1 acts as a tumour suppressor
\cite{papagiannakopoulos2023bmal1}:
its loss makes the repair/differentiation
basin epigenetically inaccessible and
stabilises the cancer attractor.

\textbf{This closes the derivation
chain.} The framework began with a
cancer cell locked in the cancer
attractor (FOXA1/EZH2 ratio, Lawson
2024). It ends by identifying that
BMAL1 loss is the molecular mechanism
by which the repair basin becomes
inaccessible --- the same geometric
failure. Cancer is cellular sleep
deprivation. The chain is circular.

\subsection{Social Information as
Primary Zeitgeber}

Social cues are confirmed co-primary
zeitgebers, entraining the circadian
clock independently of light
\cite{social2023zeitgeber}. The social
$N$ field synchronises group sleep
basins into a coupled collective
attractor.

Wild chacma baboons: dominant
individuals experience less and
lower-quality sleep than subordinates
\cite{current_biology2025baboon}.
Dominant baboons sleep closer to more
group members, experience more social
$N$ field signals during the night,
and cannot fully enter the sleep basin
without losing the social landscape
position that dominance requires.
Dominance is a landscape position
whose navigational cost does not
cease at night.

\subsection{Light and Artificial Light
at Night}

Light is the primary physical zeitgeber.
ALAN is expanding at 2.2\% per year
globally \cite{royalsoc2023alan},
disrupting the primary $N$ field
for all phototactic and circadian-
entrained species simultaneously.
600+ studies document consequences
across birds, insects, marine organisms,
mammals, and plants.

ALAN is not a minor ecological
inconvenience. It is a global
disruption of the primary $N$ field
that specifies sleep basin entry for
all life. The consequences are geometric:
organisms whose $N$ field is corrupted
lose reliable access to the basins
the field was specifying.

% ═══════════════════════════════════════��══════════════════════
\section{Subsuming the Four Existing
Theories}
\label{sec:subsuming}
% ══════════════════════════════════════════════════════════════

\subsection{Why All Four Are Correct
and Incomplete}

The four dominant theories are correct
descriptions of mechanisms within T3
(consolidated bilateral deep sleep in
endothermic organisms). This topology
covers the majority of the published
comparative sleep literature because
most studied organisms are endothermic
vertebrates.

\begin{description}[itemsep=0.3em]
  \item[SHY] describes the synaptic
    component of Axis~1 and Axis~3
    cost discharge within T3.

  \item[Energy conservation] describes
    the consequence of reduced
    navigational activity during
    Axis~2 cost recovery within T3.

  \item[Glymphatic clearance] describes
    the physical mechanism by which
    the oxidative cost substrate is
    cleared during T3 via NREM-driven
    CSF pulsation.

  \item[Memory consolidation] describes
    a downstream consequence of
    selective synaptic downscaling
    within T3.
\end{description}

None generates the dimensional invariant.
None predicts that the oldest surviving
lineages sleep the least. None predicts
convergent topologies across independent
phylogenies. None predicts that sleep
topology is a geometric property of
landscape position. The dimensional
invariant contains all four as
projections of a single geometric
structure.

\subsection{What the Framework Adds}

The framework generates eleven a-priori
predictions (P-INVARIANT-1 through
P-INVARIANT-11, stated in the Dimensional
Invariant v3.0 reasoning artifact,
2026-03-18) and five closure predictions
(P-CLOSURE-1 through P-CLOSURE-5,
stated in the Derivation Chain Closure
reasoning artifact, 2026-03-18), none
of which are derivable from any of the
four existing theories, including:

\begin{itemize}[itemsep=0.2em]
  \item Sleep duration will negatively
    correlate with lineage age across
    all phyla, topology-calibrated.

  \item Topology~2 will be found
    specifically and exclusively in
    ectothermic Navigators.

  \item T6 quiescence duration will
    scale with basin transition
    magnitude, not ongoing navigational
    demand.

  \item T8 post-emergence rebound will
    be predicted by hibernation duration
    alone.

  \item BMAL1 restoration in cancer
    cells will partially restore
    differentiation basin chromatin
    accessibility.

  \item ALAN exposure will produce
    measurable $G$ degradation
    proportional to duration and
    intensity.
\end{itemize}

% ══════════════════════════════════════════════════════════════
\section{Convergence Across Independent
Phyla}
\label{sec:convergence}
% ══════════════════════════════════════════════════════════════

\subsection{The Lock Converges Across
Seven Independent Phyla}

Seven independent phyla have arrived
at the same geometric strategy ---
deep single basin, minimal metabolism,
minimal sleep --- across evolutionary
timescales of 200--600\,Myr. The
convergence is not taxonomic. It is
geometric: the Lock strategy is the
deepest available position in $G$ for
organisms whose ecological niche is
stable and non-navigational.

The transition from T1 to T2 within
the Lock as neural complexity increases
is itself a geometric prediction: added
neural architecture adds synaptic cost
to discharge and ROS from neural firing,
but the basin remains deep enough to
keep total cost minimal. Topology shifts
but duration remains in the Lock range.

\subsection{The Navigator Converges
Across Multiple Topologies}

The Navigator strategy has independently
converged across T2, T3, and T4 in at
least seven lineages. The strategy is
the same --- ecological generalism as
Basin~$\Omega$ of ecological flexibility
--- expressed through the only topologically
feasible discharge architecture given
each lineage's full constraint set.

The octopus confirmation is the
strongest single convergence in the
dataset: 500\,Myr of separate evolution,
completely different neural architecture
(no cerebral cortex, distributed
ganglionic organisation), confirmed
Navigator ecological strategy, and
independently evolved two-phase sleep
(quiet plus active REM-like)
\cite{medeiros2021sleeping}.
Topology is determined by navigational
demand. Biology is not.

\subsection{The Unstable Zone Does Not
Converge}

No convergent feature exists in the
unstable zone except instability.
This non-convergence confirms the
framework: there is no deep basin
in the unstable zone to converge on.
Extinction is the non-convergent
attractor.

% ══════════════════════════════════════════════════════════════
\section{The Three-Channel Sonar
Instrument}
\label{sec:sonar}
% ══════════════════════════════════════════════════════════════

The complete sonar instrument for
landscape position determination has
three primary channels and two
supplementary diagnostics:

\textbf{Channel 1 --- Duration:}
Measures total discharge cost. Valid
proxy within a single topology class.
Ambiguous across topology classes.

\textbf{Channel 2 --- Topology:}
Identifies which of the eight discharge
architectures the organism uses. Removes
topology-class ambiguity from duration.
A crocodilian (T2, 2--6h true sleep)
and an elephant (T3, 2h true sleep)
have similar true sleep durations but
completely different cost structures.
Topology disambiguates.

\textbf{Channel 3 --- Behavioural
flexibility:}
Discriminates Navigator from Unstable
Zone when both occupy T3. Discriminates
T5 (flexibility impaired under pressure)
from basin restructuring (flexibility
maintained within new basin).

\textbf{Supplementary 1 --- Rebound
response:}
T5 diagnostic. T8 diagnostic (distinct
from T5: debt-to-activity ratio is
infinite in T8). Not required once
topology is identified but required
when topology data are unavailable.

\textbf{Supplementary 2 --- ROS rate:}
Predicted to discriminate ectothermic
T2 from endothermic T3 at equivalent
navigational cost levels. Testable via
metabolic rate measurement during
active period.

The migratory white-crowned sparrow
provides the critical validation of
the sonar instrument: minimal sleep
during migration with no rebound is
basin restructuring (lower actual
demand in the migratory attractor),
not T5 debt suppression. The performance
diagnostic --- maintained within-basin
performance, not impaired --- confirms
this. Duration alone would misclassify
the sparrow as approaching Lock range.
The full three-channel sonar correctly
identifies basin restructuring.

% ══════════════════════════════════════════════════════════════
\section{Discussion: The Scale
Invariance of the Framework}
\label{sec:discussion}
% ══════════════════════════════════════════════════════════════

\subsection{From Bacterium to Biosphere}

The framework operates identically
across eleven orders of magnitude of
biological scale:

\begin{description}[itemsep=0.3em]
  \item[Single-celled organisms]
    Bacterial persister cells
    demonstrate Lock/Navigator portfolio
    strategy at sub-organismal population
    level. Specialist bacteria show
    Lock-equivalent metabolic quiescence.
    Generalist bacteria show moderate
    distributed rest. The dimensional
    invariant is detectable below the
    level of the nervous system.

  \item[Developmental organisms]
    T6 topology: \textit{C.~elegans}
    lethargus is mandatory sleep at
    developmental basin transitions.
    Adolescent synaptic pruning is
    T6 operating as a prolonged landscape
    construction process --- SWA
    declining across adolescence is
    the geometric signature of a
    landscape approaching completion.

  \item[Individual organisms]
    The five primary sleep functions
    (discharge of Axes~1, 2, and 3 costs,
    somatic $G$ maintenance via growth
    hormone-driven stem cell renewal,
    and $G$ restoration via epigenetic
    maintenance) are all expressed
    within the same sleep window.

  \item[Social groups]
    Social information is a primary
    $N$ field. Group sleep is a coupled
    attractor. Dominance rank imposes
    navigational cost that does not
    cease at night. The baboon hierarchy
    is a landscape whose navigational
    demands are expressed in sleep
    quality.

  \item[Species across deep time]
    The dimensional invariant.
    The oldest surviving lineages
    sleep the least. Seven Lock phyla.
    Seven Navigator lineages across
    three topologies. The unstable zone
    does not converge. The geometry
    is consistent across 600\,Myr.

  \item[The biosphere]
    ALAN is a global primary $N$ field
    disruption affecting all species
    simultaneously. It is remodelling
    the $N$ field component of the
    $S + N + G \rightarrow R$ equation
    for thousands of species at once,
    with measurable eco-evolutionary
    consequences already underway.
\end{description}

\subsection{What Remains Open}

The framework does not yet predict
the topological consequences of a basin
that deepens in real time at
inter-generational scale (the human
Basin~B*). It does not specify the
threshold at which $G$ degradation
becomes irreversible for any given
organism. It does not model the
dynamics of the ALAN eco-evolutionary
feedback --- whether species will
adapt their $S$ substrate to function
with a corrupted $N$ field before the
disruption causes population collapse.

These are the honest limits of the
current derivation.

\subsection{The Framework is Not
Reductive}

The existing four sleep function
theories are not wrong. They are not
superseded. They are confirmed as
correct descriptions of mechanisms
within their scope (T3, endothermic
organisms). The framework does not
replace them. It provides the geometric
structure that explains why they are
all correct, why none of them alone
is sufficient, and what they are
all projections of.

The same relationship holds at every
scale. The molecular biology of BMAL1
is not superseded by the geometric
framework. It is the physical
implementation of the $N$ field at the
cellular level. The glymphatic anatomy
is not superseded. It is the physical
implementation of T3 oxidative
discharge. The Waddington epigenetic
landscape is not superseded. It is $G$.

The geometry is the structure.
The mechanisms are the geometry
made physical.

% ══════════════════════════════════════════════════════════════
\section{Conclusion}
\label{sec:conclusion}
% ══════════════════════════════════════════════════════════════

We have presented a unified geometric
theory of sleep derived from the
attractor landscape framework
$S + N + G \rightarrow R$, in which:

\begin{itemize}[itemsep=0.2em]
  \item $G$ is physically the
    epigenetic state
  \item $N$ is physically the circadian
    clock, social information, and
    light --- each a co-equal zeitgeber
  \item $S$ is the structural substrate
    maintained by sleep-dependent growth
    hormone release
  \item Sleep is the process that
    simultaneously maintains all three
    components within a single
    consolidated window
\end{itemize}

The framework makes the dimensional
invariant: across all life, at all
scales, across all of deep time, the
oldest surviving lineages sleep the
least. The minimum of the sleep spectrum
is the maximum of the survival spectrum.

The framework identifies eight sleep
topologies as geometric predictions
of landscape position, not empirical
categories.

The framework closes the derivation
chain: BMAL1 loss produces a cancer
attractor by making the repair basin
epigenetically inaccessible. Cancer
is cellular sleep deprivation. The
geometry is the same object described
from the molecular to the
civilisational scale.

Sleep is not a behaviour. Sleep is
the daily maintenance of the landscape
that makes all behaviour possible.
Survival is the depth of that
landscape.

\noindent\textcolor{rulecolor}
{\rule{\linewidth}{0.6pt}}

% ── REFERENCES ────────────────────────────────────────────────
\section*{Key References}
\label{sec:refs}

{\small
\begin{description}[itemsep=0.3em,
  labelwidth=0pt,leftmargin=1em,
  labelsep=0.5em]

  \item[\textit{Framework}]
    Lawson, E.R. (2024). The attractor
    geometry framework ($S + N + G
    \rightarrow R$) and the domestication
    syndrome as field-specified NCC
    migration arrest.
    \textit{Zenodo}.
    \href{https://doi.org/10.5281/zenodo.19094935}
    {\texttt{doi:10.5281/zenodo.19094935}}

  \item[\textit{Dimensional Invariant
    v1.0--v3.0}]
    Lawson, E.R. (2026). The Dimensional
    Invariant: Sleep Duration, Sleep
    Topology, Oxidative Cost, Attractor
    Basin Depth, and Long-Term Survival.
    \textit{Zenodo}.
    \href{https://doi.org/10.5281/zenodo.19101495}
    {\texttt{doi:10.5281/zenodo.19101495}}

  \item[\textit{Cancer (Step 1)}]
    Lawson, E.R. (2024). FOXA1/EZH2
    ratio as attractor position marker.
    \textit{Zenodo}.
    \href{https://doi.org/10.5281/zenodo.18883922}
    {\texttt{doi:10.5281/zenodo.18883922}}

  \item[\textit{SHY}]
    Tononi, G., \& Cirelli, C. (2006).
    Sleep function and synaptic
    homeostasis.
    \textit{Sleep Medicine Reviews},
    10(1), 49--62.

  \item[\textit{Glymphatic}]
    Xie, L., et al. (2013). Sleep drives
    metabolite clearance from the adult
    brain.
    \textit{Science}, 342(6156),
    373--377.

  \item[\textit{ROS death mechanism}]
    Vaccaro, A., et al. (2020). Sleep
    loss can cause death through
    accumulation of reactive oxygen
    species in the gut.
    \textit{Cell}, 181(6), 1307--1328.

  \item[\textit{Crocodilian sleep}]
    Kelly, M.L., \& Lesku, J.A. (2015).
    Bilateral slow-wave sleep in
    juvenile crocodilians.
    \textit{Journal of Experimental
    Biology}, 218, 3175--3180.

  \item[\textit{Octopus sleep}]
    Medeiros, S.L.S., et al. (2021).
    Cyclic alternation of quiet and
    active sleep states in the
    octopus.
    \textit{iScience}, 24(4), 102223.

  \item[\textit{C. elegans lethargus}]
    Allada, R., \& Siegel, J.M. (2008).
    Unearthing the phylogenetic roots
    of sleep.
    \textit{Current Biology}, 18(15),
    R670--R679.

  \item[\textit{BMAL1 tumour suppressor}]
    Papagiannakopoulos, T., et al.
    (2023). Circadian clock gene BMAL1
    acts as tumour suppressor.
    \textit{Molecular Neurobiology}.
    \href{https://doi.org/10.1007/s12035-023-03895-7}
    {doi:10.1007/s12035-023-03895-7}

  \item[\textit{Baboon dominance sleep}]
    Carter, A., et al. (2025).
    Dominant baboons experience more
    interrupted and less rest at night.
    \textit{Current Biology}.
    \href{https://www.cell.com/current-biology/fulltext/S0960-9822(25)01336-3}
    {doi:10.1016/j.cub.2025}

  \item[\textit{ALAN global disruption}]
    Gaston, K.J., et al. (2023).
    Artificial light at night: a global
    disruptor of the night-time
    environment.
    \textit{Philosophical Transactions
    of the Royal Society B}, 378,
    20220352.

  \item[\textit{Comparative sleep}]
    Melnattur, K., et al. (2023).
    Comparative biology of sleep in
    diverse animals.
    \textit{Journal of Experimental
    Biology}, 226(14), jeb245677.

  \item[\textit{Social zeitgeber}]
    Roenneberg, T., \& Merrow, M.
    (2016). The circadian clock and
    human health.
    \textit{Current Biology}, 26(10),
    R432--R443.

  \item[\textit{Glymphatic topology}]
    Nedergaard, M., et al. (2024).
    Glymphatic system review.
    \textit{Nature Reviews Neuroscience}.

\end{description}
}

% ── DOCUMENT METADATA ─────────────────────────────────────────
\vspace{0.5em}
\noindent\textcolor{rulecolor}
{\rule{\linewidth}{0.8pt}}

{\footnotesize
\textbf{Manuscript ID:}
ATTRACTOR\_GEOMETRY\_SLEEP\_THEORY\_
PAPER\_v1.0\\
\textbf{Date:} 2026-03-18\\
\textbf{Author:} Eric Robert Lawson /
OrganismCore\\
\textbf{Status:} Theory paper.
Self-contained. Cites and supersedes
the reasoning artifact series as
the primary theoretical statement.\\
\textbf{ORCID:}
\href{https://orcid.org/0009-0002-0414-6544}
{0009-0002-0414-6544}\\
\textbf{Contact:}
\href{mailto:OrganismCore@proton.me}
{OrganismCore@proton.me}\\
\textbf{Repository:}
\href{https://github.com/Eric-Robert-Lawson/attractor-oncology}
{\texttt{github.com/Eric-Robert-Lawson/
attractor-oncology}}\\
\textbf{License:} CC BY 4.0
}

\end{document}
